\section{Đánh giá kết quả đạt được}

\subsection{Kết quả đạt được so với kế hoạch}
Trong thời hạn 6 tuần, nhóm đã hoàn thành các module theo đúng 6 sprint
đã lập ở Chương 2:

\begin{table}[htbp]
    \centering
    \small
    \caption{Đối chiếu kế hoạch và kết quả thực hiện}
    \begin{tabular}{|p{4cm}|p{2cm}|p{7cm}|}
        \hline
        \textbf{Module} & \textbf{Trạng thái} & \textbf{Ghi chú} \\
        \hline
        Xác thực \& phân quyền (RBAC) & Hoàn thành & \\
        \hline
        Quản lý lịch trình nhiều ngày \& cảng ghé thăm & Hoàn thành & \\
        \hline
        Quản lý hoạt động trên tàu \& tham quan bờ & Hoàn thành & \\
        \hline
        Check-in QR/RFID/NFC & Hoàn thành & \\
        \hline
        Giao dịch POS \& đồng bộ ngoại tuyến & Hoàn thành & Module trọng
        tâm, được ưu tiên nguồn lực nhiều nhất \\
        \hline
        Tài khoản trên tàu, đối soát, hóa đơn cuối chuyến & Hoàn thành & \\
        \hline
        Dashboard vận hành & Hoàn thành & \\
        \hline
        Module BLE định vị hành khách & Chưa thực hiện & Đã xác định là
        tính năng tùy chọn ngay từ đầu, dời sang giai đoạn phát triển
        tiếp theo do giới hạn thời gian 6 tuần \\
        \hline
    \end{tabular}
\end{table}

\subsection{Đánh giá theo yêu cầu phi chức năng}
Đối chiếu với các chỉ tiêu phi chức năng đã đặt ra ở Chương 3 (thời gian
phản hồi dưới 2 giây, hỗ trợ 1.000 giao dịch/phút...), nhóm đã kiểm thử ở
quy mô nhỏ (mô phỏng tối đa vài trăm giao dịch đồng thời do giới hạn hạ
tầng thử nghiệm miễn phí) và đạt kết quả phù hợp với mục tiêu đề ra ở quy
mô thử nghiệm này; cần kiểm thử tải ở quy mô lớn hơn trước khi đưa vào
vận hành thực tế.

\subsection{Bài học kinh nghiệm}
\begin{itemize}
    \item Việc chốt sớm tech stack cụ thể (thay vì để tùy chọn
    Java/.NET/NodeJS) ngay từ Sprint 1 giúp nhóm tiết kiệm đáng kể thời
    gian trong điều kiện 6 tuần.
    \item Kiến trúc offline-first cho module POS là phần tốn nhiều công
    sức nhất nhưng cũng là điểm khác biệt cốt lõi của hệ thống so với
    các giải pháp quản lý du thuyền truyền thống đã phân tích ở Chương
    1.
    \item Thời gian dành cho kiểm thử (chỉ 1 sprint/6 sprint) là hạn chế
    lớn nhất của lịch trình rút gọn — nhóm khuyến nghị các đợt phát
    triển tiếp theo cần phân bổ tối thiểu 20\% thời gian cho kiểm thử.
\end{itemize}